\slajdid{09_1-s052}
\begin{frame}[label=09_1-s052]
\gusttekst
Ovaj iznenađujući rezultat pokazuje da suštinski u ovoj oblasti ne postoji tačka sa pojačanjem koje je jednako -1, ali isto tako i da je pojačanje konstantno (izraz važi za bilo koje $V_O$)\par\medskip
Da se vratimo na polaznu jednačinu:
\[k_{n,L}(V_{DD}-V_O-V_{Tn,L})^2=k_{n,D}(V_I-V_{Tn,D})^2\]
Mada smo rekli da to nećemo raditi ali ovde na lak način možemo izraziti $V_O$ u funkciji $V_I$
\[k_{n,L}(V_{DD}-V_O-V_{Tn,L})^2=k_{n,D}(V_I-V_{Tn,D})^2\]
\[V_O=V_{DD}-V_{Tn,L}-\sqrt{\frac{k_{n,D}}{k_{n,L}}}(V_I-V_{Tn,D})\]
pa je pojačanje u toj oblasti $\displaystyle a=\frac{dV_O}{dV_I}=-\sqrt{\frac{k_{n,D}}{k_{n,L}}}$ konstantno. I da bi dobili logičko kolo
\[\sqrt{\frac{k_{n,D}}{k_{n,L}}}>1\]
\[k_{n,D}>k_{n,L}\]
\end{frame}
